\documentclass[spanish,11pt]{article}

\usepackage[T1]{fontenc}
\usepackage[utf8]{inputenc}
\usepackage{lmodern}
\usepackage[spanish,es-nodecimaldot]{babel}
\usepackage[top=2.1cm,bottom=2.1cm,left=2.4cm,right=2.4cm]{geometry}
\usepackage{amsmath,amssymb}
\usepackage{microtype}
\usepackage{parskip}
\usepackage{xcolor}
\usepackage{hyperref}

\setcounter{secnumdepth}{5}
\setlength{\emergencystretch}{3em}

\hypersetup{
  pdftitle={Potencia de cómputo y eficiencia energética: reseña crítica de Lugo y Pedraza},
  pdfauthor={Rafael E. Marulanda},
  pdflang={es},
  colorlinks=true,
  linkcolor=blue,
  citecolor=blue,
  urlcolor=blue
}

\title{Potencia de cómputo y eficiencia energética: reseña crítica de Lugo y Pedraza}
\author{Andrés Arias, Catalina Jaramillo y Rafael Marulanda\\
  {\small Pontificia Universidad Javeriana}}
\date{4 de septiembre de 2026}

\begin{document}
\maketitle

\section{Introducción}\label{introduccion}

En su artículo \emph{Performance Evaluation for the Hash Generation
Phase of a Democratic Blockchain} (2020), Luis Lugo y Cesar Pedraza
estudian cómo distintas técnicas de computación paralela modifican el
tiempo de ejecución y el consumo energético de la generación de hashes.
El trabajo parte de un protocolo de cadena de bloques democrática
propuesto por Paul, Sarkar y Mukherjee (2014), cuyo objetivo es evitar
que la probabilidad de crear bloques dependa exclusivamente de poseer
mayor capacidad de cómputo.

En la prueba de trabajo de Bitcoin, los mineros compiten por encontrar
un hash que satisfaga una condición de dificultad. Esta competencia
convierte la potencia computacional en una variable de seguridad, pero
también en una fuente de concentración económica y demanda eléctrica.
El protocolo democrático altera la regla de selección: durante un
intervalo definido, cada nodo genera hashes y el bloque se asigna al
nodo que encuentre el valor mínimo. El proceso se divide en generación,
difusión y validación de hashes.

Lugo y Pedraza se concentran únicamente en la primera fase. Comparan
POSIX y OpenMP en CPU, CUDA y OpenCL en GPU, y Open MPI en un clúster.
Su principal resultado es que una GPU reduce el tiempo de cómputo desde
aproximadamente 1.351 segundos en la mejor configuración de CPU hasta
106--120 segundos. La aceleración se acerca a diez veces y la
implementación CUDA presenta el menor consumo estimado: 4,66 Wh frente
a cerca de 21 Wh en POSIX y OpenMP, y 100,20 Wh en Open MPI.

El artículo demuestra de forma convincente que la arquitectura de
hardware y la estrategia de paralelización importan para la eficiencia
de la generación de hashes. Sin embargo, sus resultados no bastan para
concluir que una cadena de bloques democrática sea más segura,
escalable o ambientalmente sostenible. La evaluación cubre una sola
fase, utiliza un conjunto reducido de equipos y estima la energía a
partir de especificaciones nominales, sin medir el consumo eléctrico
directamente.

\section{Contexto histórico}\label{contexto-historico}

Bitcoin integró firmas digitales, una red entre pares, bloques
encadenados mediante hashes y prueba de trabajo para resolver el doble
gasto sin una autoridad central (Nakamoto, 2008). En este diseño, la
cadena aceptada es aquella que acumula mayor trabajo computacional. La
seguridad depende de que ningún actor controle la mayoría de la
capacidad de cómputo, pero la competencia premia las economías de escala
y ha impulsado la especialización desde CPU y GPU hasta FPGA y ASIC.

Paul, Sarkar y Mukherjee (2014) propusieron una minería más democrática
para reducir esa ventaja acumulativa. En vez de premiar al primer nodo
que encuentre un hash inferior a un objetivo, su protocolo compara los
hashes generados dentro de una ventana temporal y selecciona el menor.
La regla busca ampliar la participación, mantener una tasa fija de
creación de bloques y reducir el desperdicio de cómputo. Su costo es una
mayor complejidad de mensajes: todos los nodos deben difundir resultados
y la red puede experimentar inundación de comunicaciones.

Lugo y Pedraza actualizan además el encabezado del bloque usado por ese
protocolo. Sustituyen el hash SHA-1 de 20 bytes de la dirección del
usuario por SHA-256 y amplían el campo a 32 bytes, en respuesta a la
colisión práctica de SHA-1 publicada en 2017. Para el experimento toman
datos del bloque de Bitcoin número 453194 y recorren el nonce desde cero
hasta 4003888472. El cálculo utiliza doble SHA-256 y una descomposición
cíclica que distribuye las iteraciones entre los hilos disponibles.

La comparación se ejecuta sobre un Intel Core i7-7700HQ con cuatro
núcleos físicos, una Nvidia GeForce GTX 1050 con 640 núcleos CUDA y un
clúster de cuatro instancias virtuales de Google Cloud. Este diseño
permite contrastar memoria compartida, procesamiento gráfico y memoria
distribuida bajo una carga común, aunque no representa la infraestructura
ASIC que ya dominaba la minería de Bitcoin cuando se publicó el artículo.

\section{Argumentos a favor}\label{argumentos-a-favor}

Los autores presentan un análisis de la fase de generación de hashes.
La descripción del encabezado, el rango del nonce y la plataforma experimental facilita la interpretación y la reproducción del ejercicio.

Los paradigmas POSIX y OpenMP obtienen resultados semejantes porque ambos 
utilizan mecanismos de hilos y memoria compartida equivalentes; OpenMP aporta principalmente
facilidad de programación. El rendimiento deja de mejorar cuando el
número de hilos supera la capacidad física de la CPU. En el clúster,
Open MPI mejora al pasar de uno a cuatro hilos por instancia, pero la
virtualización y la menor velocidad de los procesadores impiden que
supere a la máquina local.

El mejor resultado aparece en las GPU. Con CUDA, el tiempo cae de
1.351 segundos en la mejor configuración de CPU a 120 segundos usando
28 bloques de 128 hilos. Con OpenCL llega a 106 segundos con 24 grupos
de trabajo de 128 hilos. Los autores reportan aceleraciones de 11,28 y
12,74, respectivamente. La coherencia entre CUDA y OpenCL respalda la
conclusión de que la generación masiva de hashes se adapta bien a una
arquitectura con numerosos núcleos paralelos.

El análisis energético, aunque aproximado, agrega una dimensión que los
estudios computacionales a veces no tienen en cuenta. Los valores medios
reportados son 21,05 Wh para POSIX, 21,10 Wh para OpenMP, 100,20 Wh para
Open MPI y 4,66 Wh para CUDA. Así, el dispositivo que termina antes
también consume menos energía para completar la carga. El hallazgo
permite distinguir entre potencia instantánea y energía total: un equipo
de mayor potencia puede ser más eficiente si reduce suficientemente el
tiempo de ejecución.

El trabajo vincula dos problemas. Una regla de consenso distribuye 
oportunidades entre nodos, mientras la implementación física determina el costo de participar.
La igualdad del protocolo pierde relevancia si ejecutar el
algoritmo exige infraestructura inaccesible; por eso es pertinente medir
simultáneamente tiempo de cómputo y energía.

\section{Argumentos en contra}\label{argumentos-en-contra}

La principal limitación es el alcance. El protocolo consta
de generación, difusión y validación, pero el experimento estudia solo
la primera fase. No mide la inundación de mensajes reconocida por los
propios autores, la latencia entre nodos, el volumen de datos, la tasa de
bloques huérfanos ni el costo de llegar a consenso. Una mejora local en
la generación de hashes podría ser anulada por mayores costos de red.

Tampoco se evalúa si la regla es realmente democrática bajo condiciones
adversariales. El artículo no modela incentivos económicos, formación de
grupos, identidades múltiples, manipulación temporal ni estrategias para
ocultar hashes antes de su difusión. La igualdad de oportunidades no se
deduce únicamente de reemplazar la regla del primer hash válido por la
del hash mínimo dentro de un intervalo. Sería necesario demostrar que el
protocolo conserva seguridad y participación bajo comportamientos
estratégicos.

La medición energética es indirecta. El cálculo combina el tiempo de
ejecución, la utilización supuesta del procesador, la potencia indicada
en las fichas técnicas y una eficiencia de fuente del 85\%. No se usan
medidores en la toma eléctrica ni telemetría detallada de CPU y GPU. La
estimación tampoco incluye memoria, almacenamiento, red, refrigeración
o infraestructura del centro de datos. Además, OpenCL obtiene el menor
tiempo, pero queda fuera de la comparación energética porque la
distribución de sus grupos de trabajo es automática.

La validez externa también es limitada. El experimento utiliza una sola
CPU portátil, una GPU de consumo y cuatro máquinas virtuales, sin
intervalos de confianza ni repeticiones reportadas que permitan separar
el rendimiento medio de la variabilidad. Tampoco incluye ASIC, pese a
que estos dispositivos son la referencia económica relevante para la
minería de Bitcoin. Por ello, los resultados comparan implementaciones
de propósito general, no las alternativas realmente competitivas en esa
red.

La carga experimental reproduce el encabezado y el nonce final de un
bloque conocido. Esto ofrece una prueba controlada, pero no reproduce la
competencia de una red viva, donde cambian la dificultad, el número de
participantes, la conectividad y el costo de oportunidad de la
electricidad. El consumo de una ejecución aislada no permite inferir la
demanda total de un sistema cuyo gasto agregado responde a recompensas,
precios y entrada de nuevos mineros.

Por último, eficiencia computacional y sostenibilidad no son sinónimos.
Reducir los joules necesarios por hash puede abaratar la participación y
estimular una mayor cantidad total de hashes. Para establecer un beneficio
ambiental se necesitaría medir el consumo completo de la red, la fuente
marginal de electricidad, las emisiones asociadas y el valor social del
servicio producido. El artículo identifica correctamente un componente
del problema, pero no ofrece todavía esa evaluación económica y ambiental.

\section{Conclusión}\label{conclusion}

Lugo y Pedraza muestran que la generación de hashes es altamente sensible
a la arquitectura computacional. En su banco de pruebas, CUDA y OpenCL
reducen el tiempo alrededor de un orden de magnitud frente a la CPU, y
CUDA alcanza la mejor eficiencia energética estimada. El cambio de SHA-1
a SHA-256 para la dirección del usuario también corrige una debilidad
criptográfica del protocolo democrático original.

El aporte es un estudio de ingeniería de desempeño. Quedan
sin probar la seguridad estratégica, la escalabilidad de las comunicaciones,
el comportamiento económico de los participantes y el consumo de las
otras fases. La ausencia de medición eléctrica directa y de una comparación
con ASIC reduce además la capacidad de trasladar los resultados a Bitcoin.

Una extensión convincente debería implementar el protocolo completo en
una red con nodos heterogéneos, medir energía en la toma, incluir costos
de comunicación y refrigeración, comparar CPU, GPU y ASIC, y someter las
reglas a ataques y decisiones estratégicas. También debería separar la
eficiencia privada del equipo de la demanda agregada inducida por los
incentivos del protocolo.

Para una evaluación económica de Bitcoin, el artículo define
costos de entrada, concentración, seguridad y demanda eléctrica. 
Sin embargo, la admisibilidad social requiere un marco más amplio que incorpore
emisiones marginales, flexibilidad ante la red, localización, recuperación
de calor y distribución de beneficios. La eficiencia por hash es una
condición relevante, pero no suficiente, para afirmar que una infraestructura
digital es sostenible.

\section{Referencias}\label{referencias}

\begin{itemize}
  \setlength{\itemsep}{0pt}
  \item Lugo, L. y Pedraza, C. (2020). Performance evaluation for the hash
  generation phase of a democratic blockchain. \emph{International Journal
  of Internet Technology and Secured Transactions, 10}(3), 286--303.
  \url{https://doi.org/10.1504/IJITST.2020.107076}

  \item Nakamoto, S. (2008). \emph{Bitcoin: A Peer-to-Peer Electronic Cash
  System}. \url{https://bitcoin.org/bitcoin.pdf}

  \item Paul, G., Sarkar, P. y Mukherjee, S. (2014). Towards a more democratic
  mining in bitcoins. En \emph{Information Systems Security}, 185--203.
  Springer.

  \item Stevens, M., Bursztein, E., Karpman, P., Albertini, A. y Markov, Y.
  (2017). The first collision for full SHA-1. En \emph{Advances in Cryptology
  -- CRYPTO 2017}, 570--596.

  \item Yli-Huumo, J., Ko, D., Choi, S., Park, S. y Smolander, K. (2016).
  Where is current research on blockchain technology? A systematic review.
  \emph{PLOS ONE, 11}(10), e0163477.
  \url{https://doi.org/10.1371/journal.pone.0163477}
\end{itemize}

\end{document}
